% Generado por scripts/tablas.py a partir de Practica3_pendulos_acoplados.xlsx. No editar a mano.

\begin{table}[ht]
\centering
\caption{Modo en fase: instantes de los seis primeros máximos de cada péndulo, altura de los picos y periodo $T=(t_6-t_1)/5$. Este modo solo se registró para $\ell_1$ y $\ell_2$, pues $\omega_0$ no depende de $\ell$.}
\label{tab:modo1}
\renewcommand{\arraystretch}{1.18}\setlength{\tabcolsep}{4.2pt}
\begin{tabular}{@{}l c *{6}{r} c r@{}}
\toprule
$\ell$ & Péndulo & \multicolumn{6}{c}{Tiempo de los máximos, $t_{i}$ (\si{\milli\second})} & $p$ (\si{\volt}) & $T$ (\si{\milli\second}) \\
\cmidrule(lr){3-8}
 & & $t_1$ & $t_2$ & $t_3$ & $t_4$ & $t_5$ & $t_6$ & & \\
\midrule
\multirow{2}{*}{$\ell_{1}=\SI{25}{\centi\meter}$} & 1 & \num{2592} & \num{4588} & \num{6592} & \num{8604} & \num{10644} & \num{12620} & $-$0,13 a $-$0,15 & \num{2005.6} \\
 & 2 & \num{2712} & \num{4720} & \num{6700} & \num{8692} & \num{10676} & \num{12664} & $-$0,06 a $-$0,07 & \num{1990.4} \\
\addlinespace[2pt]
\multirow{2}{*}{$\ell_{2}=\SI{35}{\centi\meter}$} & 1 & \num{1220} & \num{3216} & \num{5204} & \num{7200} & \num{9180} & \num{11176} & $-$0,16 a $-$0,17 & \num{1991.2} \\
 & 2 & \num{13248} & \num{15264} & \num{17256} & \num{19284}\rlap{$^{\dagger}$} & \num{21292} & \num{23232} & $-$0,29 & \num{1996.8} \\
\bottomrule
\end{tabular}
\par\vspace{3pt}{\footnotesize\color{grispizarra}$^{\dagger}$\,En el cuaderno de laboratorio figuraba \SI{19824}{\milli\second}, errata de transcripción (cifras intercambiadas): con ese valor los intervalos vecinos serían de \num{2568} y \SI{1468}{\milli\second}, frente a los $\approx\SI{2000}{\milli\second}$ del resto.\par}
\end{table}
